\section{Active flux: absorb machine-specific torque production}
\label{sec:active-flux-unification}

Active flux is a second established answer to the representation question
\cite{boldea2008}. In stator-vector form,
\begin{equation}
 \boldsymbol\psi_{\mathrm a}=\boldsymbol\psi_{\mathrm s}
 -L_q\vect i_{\mathrm s}.
 \label{eq:active-flux-vector}
\end{equation}
Subtracting the isotropic \(L_q\) response leaves the part of flux that creates
torque in an equivalent nonsalient model. For a linear IPMSM,
\[
 \boldsymbol\psi_{\mathrm a}
 =[\psi_{\mathrm{pm}}+(L_d-L_q)i_d]\,\boldsymbol e_d,
 \qquad M=k_T\psi_{\mathrm a}i_q.
\]
For a stator-referred DC-excited machine,
\begin{align}
 \psi_d&=L_di_d+L_{md}i_F,& \psi_q&=L_qi_q,\\
 \boxed{\psi_{\mathrm a}}&=\boxed{L_{md}i_F+(L_d-L_q)i_d},&
 \boxed{M}&=\boxed{k_T\psi_{\mathrm a}i_q}.
 \label{eq:eesm-active-flux}
\end{align}
The derivation is merely
\(M=k_T(\psi_di_q-\psi_qi_d)\); neither sensorless operation nor a
particular observer is required for the definition. The 2009 DC-excited
active-flux control work already uses this structure
\cite{boldea2009dcactive}.

\begin{figure}[tb]
\centering
\begin{tikzpicture}[scale=1.05,>=Latex,font=\small]
\coordinate (O) at (0,0);
\coordinate (A) at (3.0,1.5);
\coordinate (B) at (1.65,.35);
\draw[->,MidnightBlue,very thick] (O)--(A) node[above]{\(\boldsymbol\psi_s\)};
\draw[->,gray,very thick] (O)--(B) node[below]{\(L_q\vect i_s\)};
\draw[->,ForestGreen!70!black,very thick] (B)--(A) node[midway,right]{\(\boldsymbol\psi_a\)};
\node[draw,rounded corners,fill=ForestGreen!8,align=center] at (6.4,.8)
 {\(\boldsymbol\psi_s=L_q\vect i_s+\boldsymbol\psi_a\)\\
 remove the isotropic response\\
 retain torque-producing flux};
\end{tikzpicture}
\caption{Active-flux vector subtraction. The remainder is the flux of the
equivalent nonsalient torque model.}
\label{fig:active-flux-subtraction}
\end{figure}

\begin{figure}[tb]
\centering
\begin{tikzpicture}[>=Latex,font=\small,node distance=8mm and 10mm]
\node[draw,rounded corners,fill=MidnightBlue!8,minimum width=2.6cm] (id)
 {\(\Delta L\,i_d\)};
\node[draw,rounded corners,fill=MidnightBlue!8,below=of id,minimum width=2.6cm] (if)
 {\(L_{md}i_F\)};
\node[draw,circle,right=of id,yshift=-6mm] (sum) {\(+\)};
\node[draw,rounded corners,fill=ForestGreen!9,right=of sum,minimum width=2.7cm] (pa)
 {\(\psi_a=s_\psi\)};
\node[draw,rounded corners,fill=Orange!12,right=of pa,minimum width=2.4cm] (tor)
 {\(\times k_Ti_q\)};
\node[draw,rounded corners,right=of tor,minimum width=1.5cm] (m) {\(M\)};
\draw[->] (id)--(sum);
\draw[->] (if)--(sum);
\draw[->,thick] (sum)--(pa);
\draw[->,thick] (pa)--(tor);
\draw[->,thick] (tor)--(m);
\node[below=9mm of pa,align=center,text width=5.0cm]
 {one flux resource, two allocation mechanisms;\\equality requires consistent current referral};
\end{tikzpicture}
\caption{EESM active flux combines reluctance and field excitation before
multiplication by the quadrature torque factor.}
\label{fig:eesm-active-flux}
\end{figure}


\subsection{Exact relation to the existing \texorpdfstring{\(s_\psi\)}{s-psi} coordinate}

The earlier coordinate
\[
 s_\psi=\Delta L i_d+L_{\mathrm{exc}}i_{\mathrm{exc}}
\]
is numerically identical to \(\psi_{\mathrm a}\) only when
\(i_{\mathrm{exc}}\) is referred to the same stator turns base and
\(L_{\mathrm{exc}}=L_{md}\). With an unreferred rotor current
\(i_F=n\,i_{\mathrm{exc}}\), the coefficient is \(nL_{md}\). Leakage belongs
in the self inductances but not in the mutual excitation term; inconsistent
use of \(L_f\), \(L_{md}\), and current referral creates an apparent scaling
difference. Under saturation, the safest definition is always the vector
identity~\eqref{eq:active-flux-vector}; a constant coefficient formula is then
only a local affine approximation.

Three related coordinate sets must not be conflated:
\begin{align*}
 (\psi_{\mathrm a},i_q,r)&\quad\text{uses physical webers and amperes},\\
 (z_\psi,z_q,z_0)&\quad\text{uses loss-normalized coordinates}.
\end{align*}
The first is immediately interpretable and portable across SPM, IPM, SynRM,
and EESM models; the residual \(r\) still needs a definition. The second makes
copper loss spherical and supplies an exact loss-orthogonal null coordinate,
but its axes depend on resistance and lose their direct units. Their link is
linear in the unsaturated model:
\(s_\psi=\beta z_\psi\) up to the stator-current whitening factor already
specified in \cref{sec:whitening}. Thus
\[
 \boxed{\text{three physical currents}
 \longrightarrow\text{two multiplicative torque factors}
 +\text{one torque-null direction}.}
\]

For an SPM, \(L_d=L_q\) and active flux is the fixed PM flux, so conventional
\(i_q\) is already an ideal torque coordinate. For an IPMSM it combines PM and
reluctance torque. For an EESM it additionally combines field and direct-axis
resources. This is why the incentive to leave raw \(dq\) currents grows with
saliency, excitation freedom, saturation, and secondary objectives.
